
\documentclass[11pt,a4paper,skipsamekey]{moderncv}
\usepackage[english]{babel}



\moderncvstyle{classic}                            
\moderncvcolor{green}                            

% character encoding
\usepackage[utf8]{inputenc}                     
\usepackage{comment}

% adjust the page margins
\usepackage[scale=0.75]{geometry}

% personal data
\name{Nathan B.}{Corral}
\address{Frankenweg 26A 53225, Bonn}
\phone[mobile]{+49 160 9178 1918}               
\email{nathan.b.corral@gmail.com}                             

\begin{document}
	\recipient{To:}{Safran Electronics \& Defense \\ Friedrich-Ebert-Straße 75 \\ 51429 Bergisch-Gladbach}
	\date{\today}
	\opening{Sehr geehrte Damen und Herren,}
	\closing{Mit freundlichen Grüßen,}
	%\enclosure[Attachment]{CV}
	\makelettertitle
	
	mit großem Interesse habe ich Ihre Stellenausschreibung für die Position als Embedded Software Engineer bei Safran gelesen. Mein akademischer Hintergrund in Intelligent Systems (M.Sc.) sowie meine Erfahrung in der Entwicklung von Embedded-Systemen und Software machen mich zu einem idealen Kandidaten für diese Position. Besonders begeistern mich die Herausforderungen in der hardwarenahen Programmierung und digitalen Hochgeschwindigkeitsdatenerfassung, die meinen Fähigkeiten und Interessen entsprechen.
	
	%Während meiner Zeit bei Aqronos habe ich teamübergreifend an der Entwicklung des Visualisierungssystems für das firmeneigene LiDAR-System gearbeitet. Dabei habe ich eine REST-API auf einem BeagleBone Black implementiert, um die Hardwareeinstellungen der Motoren zu steuern. Die Benutzer konnten diese über eine Qt-basierte C++-GUI konfigurieren.
		
	Meine Embedded-Programmierung vertiefte ich bei Head Rush Technologies, wo ich die Firmware für eine kundenspezifische Leiterplatte (PCB) entwickelte. Diese ermöglichte die Geschwindigkeitsmessung von Personen an einer Kletterwand mittels eines Zahnrad-Sensors und Interrupts. Auch wenn das System weniger komplex war als die digitalen Messgeräte von Safran, war der Prototyp erfolgreich und zuverlässig genug, um Stürze automatisch zu erkennen. Nach erfolgreichen Feldtests übergab ich den Code und die Dokumentation an das Unternehmen für die Weiterentwicklung. Dabei konnte ich meine Kenntnisse in Embedded C, Firmware-Programmierung und interdisziplinärer Zusammenarbeit weiter ausbauen.
	
	Während meines Masterstudiums in Deutschland vertiefte ich mein Wissen in Sensorfusion (z. B. Kalman-Filter) und entwickelte effiziente Methoden zur Datenverarbeitung. Zudem arbeitete ich am Humanoid Robots Lab, wo ich Erfahrungen mit GitLab CI/CD, CMake Build-Systemen und der Deployment-Optimierung für autonome Systeme sammelte.
	
	
	Nach meinem Studium habe ich aktiv an einem Open-Source-Projekt mitgearbeitet und eine neue Streaming-Audio-Transkriptionsfunktion entwickelt, die zur Veröffentlichung von Version 1.4 führte. Dabei habe ich modernen C++-Code mit Fokus auf Effizienz und Skalierbarkeit geschrieben und erfolgreich auf einem Nvidia Jetson Orin NX deployed.
	
	Ich bin von der Innovationskraft und den hohen technologischen Standards von Safran beeindruckt. Besonders reizt mich die Möglichkeit, in einem interdisziplinären Team an innovativen Messgeräten zu arbeiten. Ich bin überzeugt, dass meine Erfahrung in Embedded Systems, hardwarenaher Programmierung und Echtzeitverarbeitung eine wertvolle Ergänzung für Dein Team darstellen würde.
	
	Gerne stehe ich für ein Gespräch zur Verfügung. Ich freue mich darauf, von Dir zu hören.
	
	\vspace{0.25cm}
	\makeletterclosing
	
	\begin{comment}
		mit großem Interesse habe ich Ihre Stellenausschreibung für die Position als Embedded Software Engineer bei Safran gelesen. Mein akademischer Hintergrund in Intelligent Systems (M.Sc.) sowie meine Erfahrung in der Entwicklung von Embedded-Systemen und Software machen mich zu einem idealen Kandidaten für diese Position. Besonders begeistern mich die Herausforderungen in der hardwarenahen Programmierung und digitalen Hochgeschwindigkeitsdatenerfassung, die genau meinen Fähigkeiten und Interessen entsprechen.
		
		Während meiner Zeit bei Aqronos arbeitete ich teamübergreifend an der Entwicklung des Visualisierungssystems für das LiDAR des Unternehmens. Dabei entwickelte ich eine REST-API auf einem BeagleBone Black, die es ermöglichte, die Hardwareeinstellungen über eine C++ Qt-basierte GUI zu konfigurieren.
		
		Meine Embedded-Programmierung vertiefte ich bei Head Rush Technologies, wo ich die Firmware für eine kundenspezifische Leiterplatte (PCB) entwickelte. Diese ermöglichte die Geschwindigkeitsmessung von Personen an einer Kletterwand mittels eines Zahnrad-Sensors und Interrupts. Auch wenn das System weniger komplex war als die digitalen Messgeräte von Safran, war der Prototyp erfolgreich und zuverlässig genug, um Stürze automatisch zu erkennen. Nach erfolgreichen Feldtests übergab ich den Code und die Dokumentation an das Unternehmen für die Weiterentwicklung. Dabei konnte ich meine Kenntnisse in Embedded C, Firmware-Programmierung und interdisziplinärer Zusammenarbeit weiter ausbauen.
		
		Während meines Masterstudiums in Deutschland vertiefte ich mein Wissen in Sensorfusion (z. B. Kalman-Filter) und entwickelte effiziente Methoden zur Datenverarbeitung von Messsystemen. Zudem arbeitete ich am Humanoid Robots Lab, wo ich Erfahrungen mit GitLab CI/CD, CMake Build-Systemen und der Deployment-Optimierung für autonome Systeme sammelte.
		
		Nach meinem Abschluss habe ich aktiv an einem Open-Source-Projekt mitgearbeitet. Dabei entwickelte ich eine neue Funktion für Audio-Streaming-Transkription, die zur Version 1.4 der Software führte. Der Fokus lag auf modernem C++ mit Multithreading für maximale Effizienz und Skalierbarkeit. Abschließend habe ich das System erfolgreich auf einer Nvidia Jetson Orin NX implementiert, wodurch ich meine Fähigkeiten im Bereich GPU-gestützter Software auf Linux-Embedded-Geräten unter Beweis stellen konnte.
		
		Ich bin von der Innovationskraft und den hohen technologischen Standards von Safran beeindruckt. Besonders reizt mich die Möglichkeit, in einem interdisziplinären Team an innovativen Messgeräten zu arbeiten. Ich bin überzeugt, dass meine Erfahrung in Embedded Systems, hardwarenaher Programmierung und Echtzeitverarbeitung eine wertvolle Ergänzung für Ihr Team darstellen würde.
		
		Ich würde mich freuen, meine Motivation und Qualifikationen in einem persönlichen Gespräch näher zu erläutern. Vielen Dank für Ihre Zeit und Überlegung.
	\end{comment}
	
	
\end{document}